\begin{center}
\textsc{Gauss} an \textsc{Gerling}.\quad Göttingen, 3.~Oktober 1823.
\end{center}

\par\noindent\hspace*{\parindent}\mbox{.\ .\ .\ .\ .\ .}\quad Erlauben Sie mir noch einige \mbox{Bemerkungen.\quad\mbox{.\ .\ .\ .\ .\ .}}
\par\noindent\hspace*{\parindent}1) Beim Vergleich zweier Punkte beginne ich jedes Mal mit demjenigen,
bei dem ich am meisten befürchte, dass er ausfallen könnte.
\par\clearpage
\noindent\makebox[0pt][l]{\fontsize{13}{16}\selectfont 390}\hfill{\small\textbf{BRIEFWECHSEL MIT GERLING.}}\hfill\null\par\medskip
\par\noindent\hspace*{\parindent}2) Ist der zweite Punkt ein Heliotroplicht und besteht keine sehr große Zuversicht,
dass es sichtbar bleibt, lese ich öfter ab als gewöhnlich. Wo die Gefahr
nur äußerst gering war, habe ich bisweilen mehr als 10, d.~h.\ 15,
ja sogar 20 Messungen gemacht, bevor ich wieder ablas. Dabei setzt man
allerdings stets viel aufs Spiel. Wenn sehr zu befürchten ist, dass der zweite
Punkt ausfällt, lese ich schon nach der dritten, ja nach der zweiten oder sogar bei
jeder Beobachtung ab. Dasselbe gilt, wenn der erste Punkt etwas lange
ausbleibt.
\par\noindent\hspace*{\parindent}3) Bin ich mitten in einer Beobachtungsreihe, bei der der erste Punkt
schon mit dem Fadenkreuz eingestellt ist, lasse ich gern, bevor ich die Alhidade löse, mit einem
Handfernrohr nachsehen, ob der zweite Punkt sichtbar ist. Fällt er aber aus,
nachdem die Alhidade schon gelöst ist, und erscheint nicht wieder, ist die
Messung deswegen noch nicht verloren. Ich stelle dann nötigenfalls einen dritten Punkt mit dem Fadenkreuz ein
und bekomme so gemischte Winkel. Bei meinem diesmaligen Besuch auf dem Brocken
ist dieser Fall jedoch nur einmal eingetreten, als ich 2-mal
(Ilefeld-Hohehagen) \(+\) 2-mal (Ilefeld-Hils) nahm. Solche gemischten Winkel werden im
System ebenso gut und genau ihrem Anteil entsprechend verwendet wie reine Winkel. Im Jahr 1821 hatte
ich viele gemischte Winkel, später seltener. Deshalb habe ich an den
meisten Stationen neben den Hauptrichtungen auch Hilfsrichtungen, besonders im
Jahr 1821, als ich nur einen Heliotrop hatte. Für Hilfsrichtungen wähle ich
gern spitze Türme, die nicht allzu weit entfernt sind und möglichst nicht allzu stark
von der Horizontalebene abweichen. Diese letzte Bedingung ließ sich auf dem Brocken freilich
nicht erfüllen. Huyseburg liegt \(1^{0}\,37^{\prime}\), Hüttenrode \(1^{0}\,54^{\prime}\) und eine Stange
auf dem Wurmberg, die ich 1821 aufstellen und 1823 erneuern ließ
(wenn auch an einer etwas anderen Stelle), \(2^{0}\,13^{\prime}\) unter dem Horizont.
\par
